We now state a theorem on the relationship between the adjoint of an operator and its closure.

\begin{Theorem}
\label{thm:adjoint-closedness-and-double-adjoint-closure}
Let $(\mathcal{D}(A),A)$ be an operator from $X$ to $Y$ with dense domain. Then the adjoint operator $A^\times$ is closed. Furthermore, if $X$ and $Y$ are reflexive Banach spaces, then $(\mathcal{D}(A),A)$ is closable if and only if $\mathcal{D}(A^\times)$ is dense, in which case $\overline A=A^{\times\times}$.
\end{Theorem}

\begingroup
\postdisplaypenalty=10000
\renewcommand{\endProof}{\par\nobreak\hfill$\blacksquare$\par\vspace{0.1cm}}
\begin{Proof}
To show that $A^\times$ is closed we use Lemma~\ref{lem:closed-operator-sequential-criterion}. That is, suppose there is a sequence $v_n$ in $\mathcal{D}(A^\times)\subset Y^*$ such that $v_n\to v$ for some $v\in Y^*$ and $A^\times v_n\to f$ for some $f\in X^*$. We need to show that $v(Au)=f(u)$ for every $u\in\mathcal{D}(A)$. But since convergence in $Y^*$ $(X^*)$ implies weak-* convergence in $Y^*$ $(X^*)$, we have\[v(Au)=\lim_{n\to\infty}v_n(Au)=\lim_{n\to\infty}A^\times v_n(u)=f(u)\]for any $u\in\mathcal{D}(A)$. Thus, $A^\times$ is closed.

The rest of the proof is left to the reader (Problem~\ref{ex:adjoint-closure-theorem-proof-completion}).
\end{Proof}
\endgroup

\subsubsection{The Hilbert Adjoint Operator}\label{subsec:hilbert-adjoint-operator}

In this section we consider only operators from one Hilbert space to another. In this case we define the following operator, which is closely related to the adjoint.

\begin{Definition}
\label{def:hilbert-adjoint-operator}
Let $(\mathcal{D}(A),A)$ be a densely defined operator from a Hilbert space $H_1$ to a Hilbert space $H_2$. We define $\mathcal{D}(A^*)$ to be the set of all $v\in H_2$ such that there exists $w\in H_1$ such that
\begin{equation}\label{eq:hilbert-adjoint-pairing-definition}
(v,Au)_{H_2}=(w,u)_{H_1}
\end{equation}
for all $u\in\mathcal{D}(A)$. Note that since $\mathcal{D}(A)$ is dense, $w$ is uniquely determined by $v\in\mathcal{D}(A^*)$ and \eqref{eq:hilbert-adjoint-pairing-definition}. Thus, we can define an operator $(\mathcal{D}(A^*),A^*)$ from $H_2$ to $H_1$ by
\begin{equation}\label{eq:hilbert-adjoint-action-definition}
A^*v:=w
\end{equation}
for every $v\in\mathcal{D}(A^*)$. We call $(\mathcal{D}(A^*),A^*)$ the \textbf{Hilbert adjoint} of $(\mathcal{D}(A),A)$.
\end{Definition}

The relationship between the adjoint and the Hilbert adjoint can easily be obtained by using the Riesz maps $A_{H_1}$ and $A_{H_2}$ defined in Example~\ref{ex:riesz-map-conjugate-linear-operator}:
\begin{equation}\label{eq:hilbert-and-banach-adjoint-riesz-relation}
A^*=A_{H_1}\circ A^\times\circ A_{H_2}^{-1}.
\end{equation}
(Here, $\circ$ denotes composition of the operators.) Since a Hilbert space and its dual are isometric we rarely use the dual space directly, and when studying operators on Hilbert space we almost always use the Hilbert adjoint rather than the adjoint. In fact, it is common practice to refer to Hilbert adjoint as the adjoint, ignoring the distinction entirely. We adopt this convention, though we will use distinct notation for the two types of operators.

\begin{Theorem}
\label{thm:range-complement-adjoint-nullspace}
For any densely defined operator $(\mathcal{D}(A),A)$ from $H$ to $H$, the orthogonal complement of the range is the nullspace of the adjoint:
\begin{equation}\label{eq:range-orthogonal-complement-adjoint-nullspace}
\mathcal{R}(A)^\perp=\mathcal{N}(A^*).
\end{equation}
Furthermore, if $\mathcal{R}(A)$ is closed, then
\begin{equation}\label{eq:closed-range-adjoint-nullspace-complement}
\mathcal{R}(A)=\mathcal{N}(A^*)^\perp,
\end{equation}
i.e., the equation $Ax=f$ has a solution $x$ if and only if $f\in\mathcal{N}(A^*)^\perp$.
\end{Theorem}

\begingroup
\postdisplaypenalty=10000
\renewcommand{\endProof}{\par\nobreak\hfill$\blacksquare$\par\vspace{0.1cm}}
\begin{Proof}
We first show $\mathcal{R}(A)^\perp\subset\mathcal{N}(A^*)$. Let $z\in\mathcal{R}(A)^\perp$. Then
\begin{equation}\label{eq:range-complement-in-adjoint-domain}
(z,Au)=0=(0,u)
\end{equation}
for all $u\in\mathcal{D}(A)$. Thus, $z\in\mathcal{D}(A^*)$ and $A^*z=0$; i.e., $z\in\mathcal{N}(A^*)$. In a similar fashion, we see that if $z\in\mathcal{N}(A^*)$, then for every $u\in\mathcal{D}(A)$ we have
\begin{equation}\label{eq:adjoint-nullspace-in-range-complement}
(z,Au)=(A^*z,u)=(0,u)=0.
\end{equation}
Thus, $z\in\mathcal{R}(A)^\perp$ and $\mathcal{N}(A^*)\subset\mathcal{R}(A)^\perp$.

Finally, if $\mathcal{R}(A)$ is closed, we have
\begin{equation}\label{eq:closed-range-double-orthogonal-complement}
\mathcal{R}(A)=\mathcal{R}(A)^{\perp\perp}=\mathcal{N}(A^*)^\perp.
\end{equation}
This completes the proof.
\end{Proof}
\endgroup
